\subsection{Lie derivative and Cartan's magic formula}
In this section we discuss two differential operators on the graded algebra $\Omega^\bullet(M)$ of differential forms on a manifold $M$, namely the Lie derivative $\mathcal{L}_X$ and the interior product $\iota_X$, where $X$ is a smooth vector field on $M$. The main goal of this section is to show well-definedness of $\mathcal{L}_X$ and proof the following theorem, describing the relationships between $\mathcal{L}_X, \iota_X$ and the exterior derivative $d$.

\begin{theorem}
    \label{thm:cartan_calculus}
    Let $M$ be a manifold and $X,Y \in \mathscr{X}(M)$. Then the maps
    \begin{align*}
        d &\colon \Omega^r(M) \to \Omega^{r+1}(M), \\
        \mathcal{L}_X &\colon \Omega^r(M) \to \Omega^r(M), \\
        \iota_X &\colon \Omega^r(M) \to \Omega^{r-1}(M)
    \end{align*}
    are differential operators satisfying the following Leibniz rules for $\alpha, \beta \in \Omega^\bullet(M)$:
    \begin{align*}
        d(\alpha \wedge \beta) &= d\alpha \wedge \beta + (-1)^{|\alpha|} \alpha \wedge d\beta, \\
        \mathcal{L}_X(\alpha \wedge \beta) &= \mathcal{L}_X\alpha \wedge \beta +  \alpha \wedge \mathcal{L}_X\beta, \\
        \iota_X(\alpha \wedge \beta) &= \iota_X\alpha \wedge \beta + (-1)^{|\alpha|} \alpha \wedge \iota_X\beta,
    \end{align*}
    as well as the following identities:
    \begin{enumerate}[(i)]
        \item $d^2 = 0$, \label{eq:cartan1}
        \item $d\mathcal{L}_X = \mathcal{L}_Xd$,\label{eq:cartan2}
        \item $\mathcal{L}_X = d\iota_X + \iota_Xd$,\label{eq:cartan3}
        \item $\mathcal{L}_X\mathcal{L}_Y - \mathcal{L}_Y\mathcal{L}_X = \mathcal{L}_{[X,Y]}$,\label{eq:cartan4}
        \item $\mathcal{L}_X\iota_Y - \iota_Y\mathcal{L}_X = \iota_{[X,Y]}$,\label{eq:cartan5}
        \item $\iota_X\iota_Y + \iota_Y\iota_X = 0$.\label{eq:cartan6}
    \end{enumerate}
\end{theorem}

Firstly let us restate the definition of the Lie derivative. For a smooth manifold $M$ without boundary, $\alpha \in \Omega^r(M)$ and $X \in \mathscr{X}(M)$ we define
\[
(\mathcal{L}_X \alpha)_p := \frac{d}{dt}\bigg|_{t=0} (\phi_t^\ast \alpha)_p,
\]
where $\phi_t = \phi(t,-)$ denotes the flow of $X$.

\begin{lemma}
    \label{lem:well_definedness_Lie_derivative}
    $(\mathcal{L}_X\alpha)_p$ exists for all $p\in M$ and defines a smooth $r$-form on $M$.
\end{lemma}

Before showing the lemma, let us introduce the notion of a smooth family of $r$-forms.

\begin{definition}[Smooth family of $r$-forms]
    Let $\beta_t$ be a family of $r$-forms on $M$ indexed over some open Interval $I$. We call $\beta_t$ smooth if in some (then all) local coordinates $x^1,\ldots,x^n$ we can write
    \[
    (\beta_t)_p = \sum \beta^{i_1,\ldots,i_r}(t,p)d_px_{i_1} \wedge \ldots \wedge d_px_{i_r}
    \]
    for smooth functions $\beta^{i_1,\ldots,i_r}$. If $I$ is closed we call $\beta_t$ smooth if it can be extended to a smooth family $\beta^\prime_t$ indexed over some open open interval $I^\prime \supseteq I$.
\end{definition}

\noindent Then $t \mapsto (\beta_t)_p$ is a smooth curve in the vector space $\Lambda^r(T_p^\ast M)$, such that 
\[
\left(\frac{d}{dt}\bigg|_{t=t_0} \beta_t\right)_p := \frac{d}{dt}\bigg|_{t=t_0} (\beta_t)_p
\]
is a well-defined $r$-form. Furthermore, in local coordinates we have
\[
    \frac{d}{dt}\bigg|_{t=t_0} \beta_t = \sum \frac{\partial\beta^{i_1,\ldots,i_r}}{\partial t}(t_0,-)dx_{i_1} \wedge \ldots \wedge dx_{i_r},
\]
such that this $r$-form is also smooth and moreover depends smoothly on $t_0$. So 
\[
\dot{\beta_t} := \frac{d}{dt}\beta_t
\]
defines a smooth family of $r$-forms by
\[
\left(\frac{d}{dt}\beta_t\right)_{t_0} := \frac{d}{dt}\bigg|_{t=t_0} \beta_t.
\]

Note that if $\partial M = \emptyset$ we can also directly allow for closed intervals $I$ in the definition above. However we may then always extend the family $\beta_t$ to an open interval (in fact to all of $\R$), as $\beta_t$ can equivalently be considered a smooth $r$-form $\beta$ on $M \times I \subseteq M \times \R$. Under this correspondence we have the following interesting relationship (without proof):
\[
\dot{\beta_t} = \mathcal{L}_{\frac{\partial}{\partial t}}\beta.
\]


\begin{proof}[Proof of Lemma \ref{lem:well_definedness_Lie_derivative}]
    Let $x_1,\ldots,x_n$ be local coordinates on some open neighbourhood $U \subseteq M$ of $p$. Choose a smaller neighbourhood $U^\prime \subseteq U$ of $p$ and $\epsilon > 0$, such that the flow
    \[
    \phi \colon (-\epsilon,\epsilon) \times U^\prime \to U
    \]
    of $X$ is defined. Write
    \[
        \alpha = \sum \alpha^{i_1,\ldots,i_r}dx_{i_1} \wedge \ldots \wedge dx_{i_r}.
    \]
    Then:
    \[
        (\phi_t^\ast\alpha)_p = \sum \alpha^{i_1,\ldots,i_r}(\phi_t(p)) \,d_p(x_{i_1} \circ \phi_t) \wedge \ldots \wedge d_p(x_{i_r} \circ \phi_t)
    \]
    Expanding the differentials 
    \[
        d_p(x_{i_k} \circ \phi_t) = \sum_{j=1}^n \frac{\partial(x_{i_k} \circ \phi_t)}{\partial x_j}(p) d_px_j
    \]
    we see that $\beta_t := \phi_t^\ast \alpha$ is a smooth family of $r$-forms on $U^\prime$ indexed over $t \in (-\epsilon,\epsilon)$. Thus
    \[
    (\mathcal{L}_X \alpha)_p = \left(\frac{d}{dt}\bigg|_{t=0} \beta_t\right)_p
    \]
    exists and depends smoothly on $p$.

\end{proof}

\begin{proposition}[Basic properties of the Lie derivative]
    \label{prop:lie_derivative_properties}
    Let $\alpha \in \Omega^r(M), \beta \in \Omega^s(M)$, $\lambda \in \R$ and $X \in \mathscr{X}(M)$. Then the Lie derivative satisfies the following properties:
    \begin{enumerate}[(i)]
        \item Locality: $(\mathcal{L}_X \alpha)_U = \mathcal{L}_{X|_U} \alpha|_U$ for $U \subseteq M$ open
        \item Linearity in $\alpha$: $\mathcal{L}_X(\lambda \alpha + \beta) = \lambda \mathcal{L}_X \alpha + \mathcal{L}_X \beta$
        \item Leibniz rule: $\mathcal{L}_X(\alpha \wedge \beta) = \mathcal{L}_X \alpha \wedge \beta + \alpha \wedge \mathcal{L}_X \beta$
        \item Compatibility with $d$: $\mathcal{L}_X(d\alpha) = d(\mathcal{L}_X\alpha)$
    \end{enumerate}
\end{proposition}

\begin{proof}
    (i) follows from the fact that locally both sides are given by 
    \[
        \frac{d}{dt}\bigg|_{t=0} \beta_t
    \]
    for the smooth family $\beta_t$ defined in the proof of Lemma \ref{lem:well_definedness_Lie_derivative}. Moreover, (ii) follows immediately from the definition. To see (iii) write
    \[
    \mathcal{L}_X(\alpha \wedge \beta)_p = \frac{d}{dt}\bigg|_{t=0} (\phi_t^\ast\alpha)_p \wedge (\phi_t^\ast\beta)_p.
    \] 
    We claim: For a vector space $V$ and smooth curves $\alpha_t \in \Lambda^r(V^\ast), \beta_t \in \Lambda^s(V^\ast)$ the following holds:
    \[
        \frac{d}{dt}\bigg|_{t=0}(\alpha_t\wedge\beta_t) = \left(\frac{d}{dt}\bigg|_{t=0} \alpha_t\right) \wedge \beta_0 + \alpha_0 \wedge \left(\frac{d}{dt}\bigg|_{t=0} \beta_t\right)
    \]
    W.l.o.g. let 
    \[
    \alpha_t = a(t) \alpha^\prime, \; \beta_t = b(t) \beta^\prime 
    \]
    for some $\alpha^\prime \in \Lambda^r(V^\ast), \beta^\prime \in \Lambda^s(V^\ast)$ and smooth functions $a,b$. Then the general case follows from linearity. We compute:
    \begin{align*}
        \frac{d}{dt}\bigg|_{t=0}(\alpha_t \wedge \beta_t) &= \frac{d}{dt}\bigg|_{t=0} a(t)b(t)\alpha^\prime \wedge \beta^\prime\\
        &= \left( \frac{da}{dt}(0)b(0) + a(0)\frac{db}{dt}(0) \right)\alpha^\prime \wedge \beta^\prime\\
        &=\frac{da}{dt}(0) \alpha^\prime \wedge b(0)\beta^\prime + a(0)\alpha^\prime \wedge \frac{db}{dt}(0)\beta^\prime\\
        &= \left(\frac{d}{dt}\bigg|_{t=0} \alpha_t\right) \wedge \beta_0 + \alpha_0 \wedge \left(\frac{d}{dt}\bigg|_{t=0} \beta_t\right).
    \end{align*}

    This shows the claim and thus (iii). To see (iv) let $\beta_t$ be again the smooth family of $r$-forms defined in the proof of Lemma \ref{lem:well_definedness_Lie_derivative} and write in local coordinates $x_1,\ldots,x_n$:
    \[
    \beta_t = \sum \beta^{i_1,\ldots,i_r}(t,-)dx_{i_1} \wedge \ldots \wedge dx_{i_r}.
    \]
    By linearity we may assume that all functions $\beta^{i_1,\ldots,i_r}$ except for one vanish, i.e.
    \[
    \beta_t = f(t,-)dx_{j_1} \wedge \ldots \wedge dx_{j_r}
    \]
    for some smooth function $f$ and $1 \leq j_1 < \cdots < j_r \leq n$.
    Then
    \[
    d\beta_t = \sum_{k=1}^n \frac{\partial f(t,-)}{\partial x_k} dx^k \wedge dx_{j_1} \wedge \ldots \wedge dx_{j_r}.
    \]
    Thus
    \begin{align*}
        \mathcal{L}_X(d\alpha) &= \frac{d}{dt}\bigg|_{t=0} \phi_t^\ast d\alpha = \frac{d}{dt}\bigg|_{t=0} d\beta_t \\
        &= \sum_{k=1}^n \frac{\partial^2 f}{\partial t\partial x_k}(0,-) dx_k \wedge dx_{j_1} \wedge \ldots \wedge dx_{j_r} \\
        &= d \left(\frac{\partial f}{\partial t}(0,-) dx_{j_1} \wedge \ldots \wedge dx_{j_r}\right) \\
        &= d(\frac{d}{dt}\bigg|_{t=0} \beta_t) = d(\mathcal{L}_X\alpha)
    \end{align*}

\end{proof}

For convenience, let us also restate the definition of $\iota_X$
\begin{definition}
    Let $X \in \mathscr{X}(M)$ and $\alpha \in \Omega^r(M)$. For $r > 0$ define
    \[
    \iota_X \alpha := \alpha(X,-) \in \Omega^{r-1}(M)
    \]
    For $r=0$ let $\iota_X \alpha := 0$.
\end{definition}

Clearly $\iota_X$ is linear and only depends on the local behaviour of $X$ and $\alpha$. Furthermore, it is easily seen that 
\[
\iota_X(\alpha \wedge \beta) = \iota_X\alpha \wedge \beta + (-1)^{\abs \alpha} \alpha \wedge \iota_X\beta.
\]

Thus far the Lie derivative has only been defined on manfiolds without boundary. To extend the definition of $\mathcal{L}_X\alpha$ to manifolds with boundary, first choose extensions $X^\prime$ and $\alpha^\prime$ of $X$ and $\alpha$, respectively, to the double of $M$ and then define
\[
\mathcal{L}_X \alpha := (\mathcal{L}_{X^\prime}\alpha^\prime)|_M
\]
Of course a priori it is not clear that this definition does not depend on the choice of $X^\prime$ and $\alpha^\prime$. However, the following proposition, which is one of the formulas of Theorem \ref{thm:cartan_calculus}, shows that it is in fact independent.

\begin{proposition}[Cartan's magic formula]
    Let $X \in \mathscr{X}(M)$. Then the following identity holds:
    \[\mathcal{L}_X = d\iota_X + \iota_Xd\]
\end{proposition}
\begin{proof}
    Let $\mathcal{D}_X := d\iota_X + \iota_Xd$. From the discussion above it is easy to verify:
    \[
    \mathcal{D}_X(\alpha \wedge \beta) = \mathcal{D}_X\alpha \wedge \beta + \alpha \wedge \mathcal{D}_X\beta.
    \]
    Moreover, $\mathcal{D}_X$ clearly satisfies locality, linearity and commutes with $d$. Thus we can establish 
    \[
    \mathcal{D}_X = \mathcal{L}_X
    \]
    by showing equality on $\Omega^0(M)$. So let $f \in \Omega^0(M)$, then 
    \[
    \mathcal{D}_X f = \iota_Xdf = X \cdot f
    \]
    and
    \[
    \mathcal{L}_Xf = \frac{d}{dt}\bigg|_{t=0} (f \circ \phi_t) = df(X) = X \cdot f
    \]
    This shows the proposition.
\end{proof}

Now let us proof the remaining formulas of Theorem \ref{thm:cartan_calculus}, which are shown similiarly.

\begin{proof}[Proof of Theorem \ref{thm:cartan_calculus}]
    Equation (\ref{eq:cartan1}) is only mentioned for the sake of completeness and will not be shown here. See for example \cite{steimle_geometrie} for a proof. Equations (\ref{eq:cartan2}) and (\ref{eq:cartan3}) have already been shown above. 
    
    To see (\ref{eq:cartan4}) note that again both sides of the equation are differential operators satisfying the same Leibniz rule and commute with $d$, such that it suffices to show equality on $\Omega^0(M)$. But this follows immediately from the definition of the Lie bracket.

    We argue similiarly for (\ref{eq:cartan5}), but in this equation the differential operators on each side do not commute with $d$. However they still satisfy the same Leibniz rule, so we only need to show equality on $\Omega^0(M)$ as well as on the exact one-forms. Clearly both sides vanish on $\Omega^0(M)$ (by definition of the interior product) and for $\alpha = df \in \Omega^1(M)$ we compute:
    \begin{align*}
        (\mathcal{L}_X\iota_Y - \iota_Y\mathcal{L}_X)(df) & = \mathcal{L}_X(Y \cdot f) - (\iota_Y d \mathcal{L}_X)(f) \\
        &= XYf - YXf\\
        &= [X,Y] \cdot f = \iota_{[X,Y]}(df).
    \end{align*}
    Finally, (\ref{eq:cartan6}) is just a way to rewrite
    \[
    \alpha(X,Y,-) = -\alpha(Y,X,-).
    \]
\end{proof}

As an applictaion of Theorem \ref{thm:cartan_calculus} we can show the following identity:

\begin{corollary}
    For $\alpha \in \Omega^1(M)$ and $X,Y \in \mathscr{X}(M)$ the following identity holds:
    \[
    d\alpha(X,Y) = X \cdot \alpha(Y) - Y \cdot \alpha(X) - \alpha([X,Y])
    \]
\end{corollary}
\begin{proof}
    By Theorem \ref{thm:cartan_calculus} (\ref{eq:cartan3}) we have:
    \[
    (\iota_Y\mathcal{L}_X)(\alpha) = Y \cdot \alpha(X) + d\alpha(X,Y).
    \]
    Moreover, applying (\ref{eq:cartan5}) yields:
    \[
    (\iota_Y\mathcal{L}_X)(\alpha) = X \cdot \alpha(Y) - \alpha([X,Y]).
    \]
    The corollary follows.
\end{proof}

\headline{Integrating time dependent forms}

Let $\beta_t$ be a smooth time dependent family of differential forms on a manifold $M$ defined on some interval $I$. At the beginning of the section we described what it means to take the derivative of $\beta_t$ with respect to time. Often one is also interested in the \textit{integral} with respect to time, which is defined similiarly (i.e. pointwise). For $a ,b \in I$ and $p \in M$ we set

\[
\left( \int_{a}^{b} \beta_t dt \right)_p := \int_{a}^{b} (\beta_t)_p dt,
\]
where the right side is just the integral of a smooth curve $[a,b] \to \Lambda^r(T_p^\ast M)$ (see section \ref{sec:integrating_vector_valued_functions} for a definition of such an integral). Furthermore, in local coordinates we have
\[
    \int_{a}^{b} \beta_t dt = \sum \int_{a}^{b} \beta^{i_1,\ldots,i_r}(t,-) dt \; dx_{i_1} \wedge \ldots \wedge dx_{i_r},
\]
so $\int_{a}^{b} \beta_t dt$ is a smooth $r$-form on $M$.


\subsection{Computation rules for time-dependent forms}
In the above section we have seen how to differentiate and integrate a smooth time-dependent family $\alpha_t$ of $r-$forms. To explicitly compute such a time derivative/integral one can express $\alpha_t$ in local coordinates and take the time derivative/integral of the coefficient functions (as explained above). The goal of this section is to establish a few (coordinate free) rules for computing such expressions in practice.

\vspace{.25cm}

\noindent The Lie derivative $\mathcal{L}_X\alpha$ was defined as
\[
\frac{d}{dt}\bigg|_{t=0}\phi_t^\ast\alpha,
\]
where $\phi_t$ is the flow of $X$. We might also want to differentiate this expression at some other time $t=t_0$. The following proposition shows how this can be computed in terms of $\mathcal{L}_X\alpha$. In fact, we show a more general statement, where $\phi_t$ is the flow of a time dependent vector field.

\begin{proposition}
    \label{prop:lie_derivative_formula}
    Let $M$ be a manifold without boundary, $X_t$ a time dependent vector field on $M$ and $\alpha \in \Omega^r(M).$ Denote by $\phi_t$ the flow of $X_t$ and consider the open set $M_{t_0} = \{p\in M \mid \phi_t(p) \text{ exists for } t=t_0\}$. Then
    \[
        \frac{d}{dt}\bigg|_{t=t_0}(\phi_t^\ast\alpha)_p = \phi_{t_0}^\ast(\mathcal{L}_{X_{t_0}} \alpha)_p
    \]
    for all points $p \in M_{t_0}$.
\end{proposition}
\begin{proof}
    We begin by introducing some notation. Say $X_t$ is defined for time $t = t_1$. We then write $\phi_{t_1,t}$ for the solution of the following  nonautonomous ODE
    \[
    \frac{d}{dt}\phi_{t_1,t} = X_{t+t_1}(\phi_{t_1,t}),
    \]
    i.e. the flow of the reparametrised time-dependent vector field $X_{t+t_1}$. Then we have $\phi_{0,t} = \phi_t$ and 
    \[
    \phi_{t+t_0} = \phi_{t_0,t} \circ \phi_{0,t_0} 
    \]
    wherever this equation makes sense. Then we can compute:
    \begin{align*}
        \frac{d}{dt}\bigg|_{t=t_0}(\phi_t^\ast\alpha)_p &= \frac{d}{dt}\bigg|_{t=0}(\phi_{t+t_0}^\ast\alpha)_p\\ &= \frac{d}{dt}\bigg|_{t=0}(\phi_{0,t_0}^\ast\phi_{t_0,t}^\ast\alpha)_p \\
        &= \phi_{0,t_0}^\ast \frac{d}{dt}\bigg|_{t=0} (\phi_{t_0,t}^\ast\alpha)_p
    \end{align*}
    We claim that
    \[
        \frac{d}{dt}\bigg|_{t=0} (\phi_{t_0,t}^\ast\alpha)_p = (\mathcal{L}_{X_{t_0}}\alpha)_p
    \]
    To see this, we can explicitly express both sides in local coordinates and then use (\ref{eq:cartan_calc_further_rule_2}) from Lemma \ref{lem:cartan_calculus_further_rules} and (\ref{eq:cartan_calc_further_rule_3}) from Lemma \ref{lem:cartan_calculus_time_derivative_commutes_with_diff_ops} to reduce this to the case, where $\alpha$ is a $0$-form, which follows immediately. Essentially, this is again using the trick that two differential operators are equal if they satisfy the same Leibniz rule, commute with $d$ and agree on $0$-forms. 
\end{proof}

We now state a couple of Lemmas on how to differentiate/integrate some elementary common expressions.

\begin{lemma}
    \label{lem:cartan_calculus_further_rules}
    Let $\alpha_t, \beta_t$ be smooth time-dependent families of differential forms on a manifold $M$ and $f \colon M \to N$ a smooth map. Then the following holds:
    \begin{enumerate}[(i)]
        \item $\frac{d}{d t}f^*\alpha_t = f^*\frac{d}{d t} \alpha_t$ \label{eq:cartan_calc_further_rule_1}
        \item $\frac{d}{dt}(\alpha_t\wedge \beta_t) = \frac{d}{dt}\alpha_t \wedge \beta_t + \alpha_t \wedge \frac{d}{dt}\beta_t$ \label{eq:cartan_calc_further_rule_2}
    \end{enumerate} 
\end{lemma}
\begin{proof}
    For (\ref{eq:cartan_calc_further_rule_1}) note that for a given point $p\in M$ the pullback along $f$ is given by a linear map
    \[
        f^\ast \colon\Lambda^r(T_{f(p)}^\ast N) \to \Lambda^r(T_{p}^\ast M),
    \] 
    so the statement follows directly from the chain rule. (\ref{eq:cartan_calc_further_rule_2}) follows from the ordinary Leibniz rule for smooth functions and can be seen by computing the derivative in local coordinates. We have in fact already seen this computation in the proof of Proposition \ref{prop:lie_derivative_properties}.
\end{proof}

This next Lemma shows, that the time derivative of time-dependent differential forms commutes with the differential operators from Theorem \ref{thm:cartan_calculus}.

\begin{lemma}
    \label{lem:cartan_calculus_time_derivative_commutes_with_diff_ops}
    Let $\alpha_t$ be a time-dependent family of differential forms and $X$ a vector field on a manifold $M$. Then the following holds:
    \begin{enumerate}[(i)]
        \item $\frac{d}{dt}d\alpha_t = d\frac{d}{dt}\alpha_t$ \label{eq:cartan_calc_further_rule_3}
        \item $\frac{d}{dt} \mathcal{L}_X \alpha_t = \mathcal{L}_X \frac{d}{dt}\alpha_t$ \label{eq:cartan_calc_further_rule_4}
        \item $\frac{d}{dt} \iota_X \alpha_t = \iota_X \frac{d}{dt}\alpha_t$ \label{eq:cartan_calc_further_rule_5}
    \end{enumerate}
\end{lemma}
\begin{proof}
    (\ref{eq:cartan_calc_further_rule_3}) follows from Schwarz's Theorem by computing the expression in local coordinates. This computation has essentially already been carried out in the proof of Proposition \ref{prop:lie_derivative_properties}. (\ref{eq:cartan_calc_further_rule_4}) follows from (\ref{eq:cartan_calc_further_rule_3}) and (\ref{eq:cartan_calc_further_rule_5}) using Cartan's magic formula. For (\ref{eq:cartan_calc_further_rule_5}) note that for $p \in M$ the map 
    \[
    \iota_{X_p} \colon \Lambda^r(T_{p}^\ast M) \to \Lambda^{r-1}(T_{p}^\ast M)
    \]
    is linear so the statement follows again from the chain rule.
\end{proof}

Moreover, the integral with respect to time also commutes with the differential operators, as the following Lemma shows.

\begin{lemma}
    Let $\alpha_t$ be a time-dependent family of differential forms defined on some interval $I$ and $X$ a vector field on a manifold $M$. Given $a,b \in I$ the following holds: 
    \begin{enumerate}[(i)]
        \item $d \int_{a}^{b} \alpha_t dt = \int_{a}^{b} d\alpha_t dt$ \label{eq:cartan_calc_further_rule_6}
        \item $\mathcal{L}_X \int_{a}^{b} \alpha_t dt = \int_{a}^{b} \mathcal{L}_X\alpha_t dt$\label{eq:cartan_calc_further_rule_7}
        \item $\iota_X \int_{a}^{b} \alpha_t dt = \int_{a}^{b} \iota_X\alpha_t dt$\label{eq:cartan_calc_further_rule_8}
    \end{enumerate}
\end{lemma}
\begin{proof}
    For (\ref{eq:cartan_calc_further_rule_6}) choose local coordinates $x_1,\ldots,x_n$ and let $\alpha^{i_1,\ldots,i_r}(t,x)$ be a component function of $\alpha_t$ with respect to these coordinates. Then 
    \[
    \frac{\partial}{\partial x_i} \int_{a}^{b} \alpha^{i_1,\ldots,i_r}(t,x) dt = \int_{a}^{b}  \frac{\partial}{\partial x_i} \alpha^{i_1,\ldots,i_r}(t,x) dt
    \]
    since $[a,b]$ is compact. The statement follows. Again (\ref{eq:cartan_calc_further_rule_7}) follows from (\ref{eq:cartan_calc_further_rule_6}) and (\ref{eq:cartan_calc_further_rule_8}) using Cartan's magic formula. (\ref{eq:cartan_calc_further_rule_8}) follows from linearity of the integral (see Lemma \ref{lem:linearity_of_integral}) by again considering the linear map $\iota_{X_p}$ for $p \in M$.
\end{proof}

Finally, we explain how to deal with time dependent vector fields in the differential operators.

\begin{lemma}
    Let $X_t$ be a time-dependent vector field on a manifold $M$. Then the following holds: 
    \begin{enumerate}[(i)]
        \item $\frac{d}{dt} \mathcal{L}_{X_t} = \mathcal{L}_{\frac{d}{dt}X_t}$ \label{eq:cartan_calc_further_rule_9}
        \item $\frac{d}{dt} \iota_{X_t} = \iota_{\frac{d}{dt}X_t}$ \label{eq:cartan_calc_further_rule_10}
    \end{enumerate}
\end{lemma}
\begin{proof}
    (\ref{eq:cartan_calc_further_rule_9}) follows from (\ref{eq:cartan_calc_further_rule_10}) as well as Lemma \ref{lem:cartan_calculus_time_derivative_commutes_with_diff_ops} (\ref{eq:cartan_calc_further_rule_3}) using Cartan's magic formula. For (\ref{eq:cartan_calc_further_rule_10}) let $\alpha \in \Omega^r(M)$ and consider the linear map 
    \[
    A \colon T_p M \to \Lambda^{r-1}(T_p^\ast M), \;\, v \mapsto \alpha_p(v,-)
    \]
    for $p \in M$. Then 
    \[
    \left( \frac{d}{dt} \iota_{X_t} \alpha \right)_p = \frac{d}{dt} A(X_t(p))
    \]
    and so the statement follows again from the chain rule.
\end{proof}

\headline{Moser`s trick}

In many cases one has to differentiate a time-dependent family of differential forms, which is given by an expression where the time variable appears at multiple instances. We first proof a formula for the derivative of such an expression by discussing an important example, namely Moser`s trick, which is used to proof the Darboux Theorem about symplectic/contact manfiolds. Following that we present a straightforward cooking recipe illustrated by another example.

Let $M$ be a manifold (without boundary). Given a smooth family $\alpha_t$ of differential forms on $M$ indexed over $t \in [0,1]$ we want to find a family of diffeomorphisms $\varphi_t \colon M \to M$ (i.e. an isotopy), such that $\varphi_0 = id_M$ and 
\[
    \varphi_t^\ast \alpha_t = \alpha_0 \: \text{ for all } t \in [0,1],
\]
or equivalently:
\begin{equation}
    \label{eq:mosers_trick}
    \frac{d}{dt} \varphi_t^\ast \alpha_t = 0.    
\end{equation}


% \begin{remark}
%     If $\partial M \neq \emptyset$, we call $\varphi_t$ a smooth family of diffeomorphisms if the map 
%     \[
%     M \times [0,1] \to M, \: (p,t) \mapsto \varphi_t(p)
%     \]
%     is smooth in the sense of Definition \ref{def:extension_smooth}, where we consider $M \times [0,1] \subseteq M \times \R$.
    % If $\partial M \neq \emptyset$ we say that $\varphi_t$ is a smooth family of diffeomorphisms if it can be extended to a smooth family $\varphi^\prime_t$ indexed over $\R$. In fact, if $\partial M = \emptyset$ we can also extend $\varphi_t$ to $\R$, using the following trick: The isotopy $\varphi_t$ corresponds to the level preserving diffeomorphism
    % \[
    % \Phi \colon M \times [0,1] \to M \times [0,1], (p,t) \mapsto (\varphi_t(p),t).
    % \]
    % This diffeomorphism then defines a vector field $Y:= \Phi_\ast \frac{\partial}{\partial t}$, whose integral curves uniquely determine $\varphi_t$ (this trick is also explained in \cite{Milnor1965hcobord} Thm. 5.6). We can extend this vector field arbitrarily to $M \times \R$ and then set the component pointing in the $\R$ direction to $\frac{\partial}{\partial t}$. 
%\end{remark}

Directly constructing such an isotopy is difficult but we can make use of the following trick: The isotopy $\varphi_t$ corresponds to the level preserving diffeomorphism
\[
\Phi \colon M \times [0,1] \to M \times [0,1], (p,t) \mapsto (\varphi_t(p),t).
\]
This diffeomorphism then defines a vector field $\Phi_\ast \frac{\partial}{\partial t}$ on $M \times [0,1]$, whose integral curves uniquely determine $\varphi_t$. This suggests that in order to construct an isotopy $\varphi_t$ we can construct a vector field $Y$ on $M \times [0,1]$, of the form 
\[
Y_{(p,t)} = ((X_t)_p, \frac{\partial}{\partial t}) 
\]
for a time dependent vector field $X_t$ on $M$ (this trick is also explained in \cite{Milnor1965hcobord} Thm. 5.6). 

Then (\ref{eq:mosers_trick}) can be reformulated as a condition on $X_t$, which is shown in the following Lemma.

\begin{lemma}
    \label{lem:mosers_trick_lemma}
    In the situation above the following holds:
    \[
    \frac{d}{dt} \varphi_t^\ast \alpha_t = \phi_t^\ast(L_{X_t}\alpha_t + \dot{\alpha_t})
    \]
\end{lemma}

\begin{proof}
    Assume that $\alpha_t$ is defined for all $t \in \R$. The idea is 
    to "decouple" the time variables and consider the two parameter family
    \[
        g \colon [0,1] \times \R \to T_p^\ast M, \: (t,s) \mapsto (\phi_t^\ast \alpha_s)_p
    \]
    The partial derivatives are now much easier to compute. By Proposition \ref{prop:lie_derivative_formula} we have:
    \[
    \frac{\partial}{\partial t}\bigg|_{t=t_0} \varphi_t^\ast \alpha_s = \varphi_{t_0}^\ast(\mathcal{L}_{X_{t_0}} \alpha_s)
    \]
    Moreover, by Lemma \ref{lem:cartan_calculus_further_rules}
    \[
        \frac{\partial}{\partial s}\bigg|_{s=t_0} \varphi_t^\ast \alpha_s = \varphi_t^\ast \left(\frac{\partial}{\partial s}\bigg|_{s=t_0}  \alpha_s \right) = \varphi_t^\ast \dot{\alpha}_{t_0}.
    \]
    In order to compute the derivative of 
    \[
    f \colon [0,1] \to T_p^\ast M, \: t \mapsto (\phi_t^\ast \alpha_t)_p
    \]
    consider the diagonal map
    \[
    \Delta \colon [0,1] \to [0,1] \times \R, \: t \mapsto (t,t).
    \]
    Then
    \[
    f = g \circ\Delta,
    \]
    hence
    \begin{align*}
        \frac{d}{dt} \bigg|_{t = t_0} f(t) & = D_{t_0}f(1) = D_{(t_0,t_0)}g \circ D_{t_0}\Delta (1) \\
        & = D_{(t_0,t_0)}(1,1) = D_{(t_0,t_0)}(1,0) + D_{(t_0,t_0)}(0,1) \\
        & = \frac{\partial g}{\partial t} (t_0,t_0) + \frac{\partial g}{\partial s} (t_0,t_0)  
    \end{align*}
    The Lemma follows.
\end{proof}
So (\ref{eq:mosers_trick}) can be reformulated to the condition:
\[
    0 = L_{X_t}\alpha_t + \dot{\alpha_t} = d\iota_{X_t}\alpha_t + \iota_{X_t}d\alpha_t + \dot{\alpha_t},
\]
where we used Theorem \ref{thm:cartan_calculus} (\ref{eq:cartan3}) in the second equation. Hence it suffices to construct a time dependent vector field $X_t$ satisfying the above conditon and whose integral curves exist for time $t = 1$. 

\headline{A cooking recipe}

The proof of Lemma \ref{lem:mosers_trick_lemma} gives a cooking recipe for computing derivatives of smooth families of differential forms indexed over some time $t$, where $t$ appears at multiple instances. Namely, fix $t$ for all but one instance and differentiate. Then take the sum of all of the results. We illustrate this in an example.

\begin{example}
    Let $X_t$ be some time dependent vector field on a manifold $M$ with flow $\phi_t$. Furthermore, let $\alpha_t$ be a smooth family of differential forms and $f_t$ a smooth family of functions on $M$. We want to compute:
    \[
    \frac{d}{dt}\phi_t^\ast(e^{f_t}\alpha_t).
    \]
    Begin by fixing all but the first occurence of $t$ and differentiate, i.e. compute
    \begin{align*}
        \frac{\partial}{\partial t} \phi_t^\ast(e^{f_s}\alpha_s) = \phi_t^\ast(\mathcal{L}_{X_t}(e^{f_s}\alpha_s)) &= \phi_t^\ast((X_t \cdot e^{f_s})\alpha_s + e^{f_s}\mathcal{L}_{X_t}\alpha_s)\\ &= \phi_t^\ast((X_t \cdot f_s) e^{f_s} \alpha_s + e^{f_s}\mathcal{L}_{X_t}\alpha_s)
    \end{align*}
    Here we used Proposition \ref{prop:lie_derivative_formula} and the Leibniz rule for the Lie derivative. Next fix all but the second occurence of $t$ and using Lemma \ref{lem:cartan_calculus_further_rules} compute:
    \[
        \frac{\partial}{\partial t} \phi_s^\ast(e^{f_t}\alpha_s) = \phi_s^\ast(\frac{\partial }{\partial t}e^{f_t}\alpha_s) = \phi_s^\ast(\dot{f_t}e^{f_t}\alpha_s)
    \]
    And lastly
    \[
        \frac{\partial }{\partial t} \phi_s^\ast(e^{f_s}\alpha_t) = \phi_s^\ast(e^{f_s}\frac{\partial }{\partial t} \alpha_t) = \phi_s^\ast(e^{f_s}\dot{\alpha_t})
    \]
    Finally plug in $s=t$ and sum the results:
    \begin{align*}
        \frac{d}{dt}\phi_t^\ast(e^{f_t}\alpha_t) &= \phi_t^\ast((X_t \cdot f_t) e^{f_t} \alpha_t + e^{f_t}\mathcal{L}_{X_t}\alpha_t) + \phi_t^\ast(\dot{f_t}e^{f_t}\alpha_t) + \phi_t^\ast(e^{f_t}\dot{\alpha_t})\\
        &= \phi_t^\ast(e^{f_t}((X_t \cdot f_t) \alpha_t + \mathcal{L}_{X_t}\alpha_t + \dot{f_t}\alpha_t + \dot{\alpha_t}))
    \end{align*}
\end{example}